\texttt{HAETAE}, a lattice-based Fiat--Shamir signature scheme proposed for the
post-quantum era, may be vulnerable to fault injection attacks (FIA) in real
implementations regardless of its mathematical security. In particular, the response
computation $\zvec=\yvec+(-1)^b c\svec_1$ is the only step in which the secret key and
the one-time nonce are combined; if its integrity is not guaranteed, secret information
can be exposed directly. We analyze the \texttt{HAETAE} signing procedure and identify
three fault points in the response computation (skipping $c\svec$, skipping $+\yvec$, and
rejection bypass). We show that skipping the $+\yvec$ addition enables the secret key
$\svec_1$ to be recovered directly from a \emph{single} faulted signature. On a
full-signature STM32F405 implementation with ChipWhisperer-Husky, under an
\textbf{emulated single-fault model}, we recover all 768 coefficients of $\svec_1$ with
100\% success from one faulted signature, and the recovery pipeline is confirmed by a
self-test. We further propose \texttt{IRV}
(Infective Randomized Verification), a
lightweight countermeasure that unifies recomputation-based verification, signing-boundary
norm re-checking, and secret-key integrity checking into \textbf{branchless infective
masking} that does not rely on comparison (branch) operations. Experimentally, IRV
(i)~blocks all seven fault points including rejection bypass, (ii)~resists second-order
faults that skip detection branches, and (iii)~operates at less than half the
\emph{signing time} of full double computation ($+14\%$ signing-time overhead), though it
requires more code/RAM than double computation.

\smallskip
\noindent\textbf{Keywords.} Post-quantum cryptography; Lattice-based signature;
\texttt{HAETAE}; Fault injection attack; Infective countermeasure; Single-signature key recovery.
